% Reusable exam scaffold: replace bracketed text and symbolic marks before use.
% Build twice for question references; see the build guide for TEXINPUTS.
\makeatletter
\def\input@path{{preamble/}{../preamble/}{../../preamble/}{../../../preamble/}{../../../../book2026-2027/preamble/}}
\makeatother
\input{preamble_exams}
\graphicspath{{./}{preamble/}{../preamble/}{../../preamble/}{../../../preamble/}{../../../../book2026-2027/preamble/}}

% Assessment-specific settings. Symbolic marks are NOT suggested allocations.
\setYColumnWidth{0.23}
\newcommand{\CriterionA}{[ID A]}
\newcommand{\CriterionB}{[ID B]}
% Ledger: Q1 A=a+b; Q2 B=c+d; Q3 A=e. Totals: A=a+b+e, B=c+d.
% Replace all task marks, heading sums and criterion totals consistently.

% Preserve starred/unstarred calls and subtract color-box padding ONCE.
% Do not reduce linewidth again at individual calls.
\let\ExamSubsectionBox\SubsectionBox
\RenewDocumentCommand{\SubsectionBox}{s m}{%
  \par\Needspace{4\baselineskip}% keep heading with opening prompt
  \begingroup
  \setlength{\linewidth}{\dimexpr\linewidth-2\fboxsep\relax}%
  \IfBooleanTF{#1}{\ExamSubsectionBox*{#2}}{\ExamSubsectionBox{#2}}%
  \endgroup\nobreak
}
% Shared headings use stepcounter, not refstepcounter: capture the automatic
% number explicitly for score-table references, without manual question IDs.
\makeatletter
\newcommand{\QuestionLabel}[1]{%
  \begingroup\edef\@currentlabel{\arabic{subsectionboxproblem}}%
  \label{#1}\endgroup
}
\makeatother
\newsavebox{\ScoreTableA}
\newsavebox{\ScoreTableB}
\newlength{\ScoreTableGap}
\setlength{\ScoreTableGap}{1em}
\newif\ifSeparateActivity
\SeparateActivitytrue % Optional: set false to omit the reminder.

\begin{document}
% Full-width editable metadata. Alternatively use ExamHeader when its fixed
% teacher/year match the brief. Keep all front matter outside multicols.
\noindent
\begin{tabularx}{\textwidth}{|C{2.8cm}|C{\dimexpr\textwidth-5.6cm-6\tabcolsep-4\arrayrulewidth\relax}|C{2.8cm}|}
\hline
\centering\vspace{3mm}\includegraphics[width=2.5cm]{logo.png}
& \CellCenter{\textbf{GLOBAL ECONOMICS}\par
  \textbf{Grade [grade] -- [assessment designation]}\par
  \textbf{[Exam title]}\par\textbf{[Teacher]}}
& \CellCenter{\textbf{[Term]}\par\textbf{[School year]}}\\\hline
\end{tabularx}\par\smallskip
\LearningObjective{[Learning objective: observable knowledge and skills.]}{%
\textbf{\CriterionA}: [Description of the first assessed criterion.]\par
\textbf{\CriterionB}: [Description of the second assessed criterion.]
% Add criteria as needed; distinguish separately assessed criteria.
}\par\smallskip
\StudentInfo\par

\begin{multicols}{2}
\raggedcolumns
\setlength{\parskip}{3pt}
\SubsectionBox*{Criteria assessment}
[Assessed criteria and available points per criterion; use the totals below.]
[Passing thresholds, if applicable; specify this exam's points or cumulative
points.] [Required procedures and justifications.] [Partial-credit rule.]
[Penalties, triggers and scope, if applicable.] [Answer location: this sheet
or identified separate paper; submission requirements.]
% Assessment-specific policies: remove inapplicable clauses.

\noindent\begin{minipage}{\linewidth}
\centering\small
\setlength{\tabcolsep}{4pt}\renewcommand{\arraystretch}{1.05}
\textbf{Available points}\par\smallskip
% One bordered table per criterion; include only its scored questions.
\sbox{\ScoreTableA}{\begin{tabular}{|c|c|}
\hline\textbf{Question} & \textbf{\CriterionA}\\\hline
\ref*{q:context} & $a+b$\\\hline
\ref*{q:short} & $e$\\\hline
\textbf{Total} & $a+b+e$\\\hline
\end{tabular}}%
\sbox{\ScoreTableB}{\begin{tabular}{|c|c|}
\hline\textbf{Question} & \textbf{\CriterionB}\\\hline
\ref*{q:ordered} & $c+d$\\\hline
\textbf{Total} & $c+d$\\\hline
\end{tabular}}%
% Measure the pair; automatically stack if it does not fit the column.
\ifdim\dimexpr\wd\ScoreTableA+\ScoreTableGap+\wd\ScoreTableB\relax>\linewidth
\usebox{\ScoreTableA}\par\smallskip\usebox{\ScoreTableB}%
\else
\usebox{\ScoreTableA}\hspace{\ScoreTableGap}\usebox{\ScoreTableB}%
\fi
\end{minipage}\par
% Repeat in pairs for more criteria. Each individual table must fit linewidth.
% Split tall tables by labeled question ranges; repeat column headers.
% Exclude the separate activity from every score table.

% LONG QUESTION: no enclosing minipage; permit flow across columns/pages.
\SubsectionBox{[Contextual question]\score{\CriterionA-\sml{$a+b$}}}
\QuestionLabel{q:context}
[Introduce the decision maker, context, alternatives, uncertain outcomes,
units, assumptions and parameter domains. Define all symbols used below.]

% Keep each short caption/table block together.
\noindent\begin{minipage}{\linewidth}
\small\textbf{[Payoffs in stated units]}\par\smallskip
\setlength{\tabcolsep}{3pt}
\begin{tabularx}{\linewidth}{@{}Xcc@{}}
\toprule Alternative & [State I] & [State II]\\\midrule
{[Probability, if given]} & $p$ & $1-p$\\
{[Alternative I]} & $x_{11}$ & $x_{12}$\\
{[Alternative II]} & $x_{21}$ & $x_{22}$\\\bottomrule
\end{tabularx}
\end{minipage}\par
% Delete probabilities when unavailable; replace symbols with verified data.
\noindent\begin{minipage}{\linewidth}
\small\textbf{[Decision statistics]}\par
[Define each statistic and whether larger or smaller is preferred.]
\setlength{\tabcolsep}{3pt}
\begin{tabularx}{\linewidth}{@{}Xccc@{}}
\toprule Alternative & $T_1$ & $T_2$ & $T_3$\\\midrule
{[Alternative I]} & $t_{11}$ & $t_{12}$ & $t_{13}$\\
{[Alternative II]} & $t_{21}$ & $t_{22}$ & $t_{23}$\\\bottomrule
\end{tabularx}
\end{minipage}\par
% Adapt statistics to the topic; check consistency with payoff data.
[Inline task: calculate or compare using the supplied information.]
\sscore{\CriterionA-\sml{$a$}}
[Inline task: interpret, justify or investigate a change; state domains,
eligibility and tie rules where needed.] \sscore{\CriterionA-\sml{$b$}}

% ORDERED TASKS: let longer questions break between tasks.
\SubsectionBox{[Ordered procedure]\score{\CriterionB-\sml{$c+d$}}}
\QuestionLabel{q:ordered}
[Supply self-contained context and data, or explicitly identify shared data
above. Specify the operation order and any selection or tie rules.]
\begin{enumerate}[label=(\alph*)]
\item {[Apply the first stage; show the procedure and intermediate result.]}
  \sscore{\CriterionB-\sml{$c$}}
\item {[Apply the remaining stages in order; explain the final result.]}
  \sscore{\CriterionB-\sml{$d$}}
\end{enumerate}
% For mixed criteria list each criterion sum in the heading; enter the question
% in each relevant score table with only that criterion's share.

% SHORT QUESTION: keep heading, prompt and optional workspace together.
\noindent\begin{minipage}{\linewidth}
\SubsectionBox{[Short question]\score{\CriterionA-\sml{$e$}}}
\QuestionLabel{q:short}
[Give the complete prompt and any data needed for a brief response.]
\sscore{\CriterionA-\sml{$e$}}
% For on-sheet answers add appropriately sized space, e.g. \vspace{15mm}.
\end{minipage}\par
\ifSeparateActivity
\noindent\begin{minipage}{\linewidth}
\SubsectionBox*{Separate activity reminder (unscored)}
[Activity title and separately assessed criterion.] [Required deliverable,
submission location and reference to the activity's own instructions.]
This reminder carries no points in this exam.
\end{minipage}\par
\fi
\end{multicols}
\end{document}
